\subsection{Réactivité \& Classification : Technologies au lithium}

\setcounter{equation}{0}
\setcounter{Q}{0}

\cteacher{Attention, cet exercice va poser quelques difficultés. Il faut leur montrer qu'équilibrer des équations dans un milieu \ce{Li^+}, c'est comme équilibrer un équilibre redox en milieu acide où le couple Ox/Réd fait appel à un échange de \ce{H+}.
Au vu des questions en sortie de CM, ils n'ont pas compris que les familles de réactions présentées ne sont que des catégories et que ça reste des processus redox dans tout les cas. 
}

Dans les technologies de stockage électrochimiques au lithium, trouver les demi-équations redox associées à chacun des pôles (en décharge) et proposer une équation de réaction fictive modélisant la réaction.

Vous préciserez le type de réaction mise en jeu dans ces systèmes.

\Q $\ominus$ : \ce{C_6}/\ce{LiC6}, $\oplus$ \ce{CoO2} / \ce{LiCoO2}, électrolyte contenant \ce{Li^+} en phase organique

\solution{ 
\begin{minipage}{0.48\textwidth}
\begin{align*}
\ce{LiC6 (s) &= C_6 (s) + Li^+ (org) + e^-} \\
\ce{CoO2(s) + Li^+(org) + e^- &= LiCoO2(s)} \\
\hline
\ce{1/6 LiC6(s) + CoO2(s) &= LiCoO2 (s) + C_6 (s)}
\end{align*}
\end{minipage}\begin{minipage}{0.48\textwidth}
\begin{center}
Réaction d'insertion (plus précisément, intercalation de \ce{Li^+})
\end{center}
\end{minipage}}



\Q $\ominus$ \ce{Li^+}/\ce{Li} , $\oplus$ \ce{Li_{15}Si_4}/\ce{Si}

\solution{ 
\begin{minipage}{0.4\textwidth}
\begin{align*}
\ce{Li(s) &= Li^+ + e^-} \\
\ce{Si + 15/4 Li^+ + 15/4 e^- &= 1/4 Li_{15}Si_4} \\
\hline
\ce{Si(s) + 15/4 Li(s) &= 1/4 Li_{15}Si_4 (s)}
\end{align*}
\end{minipage}\begin{minipage}{0.4\textwidth}
\begin{center}
Réaction de formation (ici un alliage Si-Li)
\end{center}
\end{minipage}}

\Q $\ominus$ \ce{Li^+}/\ce{Li} , $\oplus$ \ce{Li_{15}Sn_5}/\ce{Sn}

\solution{ 
\begin{minipage}{0.4\textwidth}
\begin{align*}
\ce{Li(s) &= Li^+ + e^-} \\
\ce{Sn(s) + 15/5 Li^+ + 15/5 e^- &= 1/5 Li_{15}Sn_5 (s)} \\
\hline
\ce{Sn(s) + 15/5 Li(s) &= 1/5 Li_{15}Sn_5 (s)}
\end{align*}
\end{minipage}\begin{minipage}{0.4\textwidth}
\begin{center}
Réaction de formation d'un alliage Sn-Li
\end{center}
\end{minipage}}

\Q $\ominus$ \ce{Li^+}/\ce{Li} , $\oplus$ \ce{CoS}/\ce{Co}

\solution{ 
\begin{minipage}{0.4\textwidth}
\begin{align*}
\ce{Li(s) &= Li^+ + e^-} \\
\ce{CoS (s) + 2 e^- + 2 Li^+ &= Co(s) + Li2S(s)} \\
\hline
\ce{Li(s) + CoS(s) &= Co(s) + Li2S(s)}
\end{align*}
\end{minipage}\begin{minipage}{0.4\textwidth}
\begin{center}
Réaction de déplacement (de "\ce{S}")
\end{center}
\end{minipage}}

\Q $\ominus$ \ce{Li^+ / Li}, $\oplus$ \ce{Cu4Mo6S8}/\ce{Cu}

\solution{ 
\begin{minipage}{0.4\textwidth}
\begin{align*}
\ce{Li(s) &= Li^+ + e^-} \\
\ce{Cu4Mo6S8(s) + 4Li^+ + 4e^- &= Li4Mo6S8(s) + 4Cu(s)} \\
\hline
\ce{Li(s) + Cu4Mo6S8(s) &= Li4Mo6S8(s) + 4Cu(s)}
\end{align*}
\end{minipage}\begin{minipage}{0.4\textwidth}
\begin{center}
Réaction de déplacement (de "\ce{Mo6S8}")
\end{center}
\end{minipage}}